\documentclass[12pt]{article}
\usepackage{float}
\usepackage{graphicx}
\usepackage{fontspec}
\usepackage{url}
\usepackage{apacite}
\setmainfont{Times New Roman}
\usepackage[letterpaper,top=2cm,bottom=2cm,left=2cm,right=2cm,marginparwidth=1.5cm]{geometry}

\title{Weekly assignment 7: Ethical Frameworks}
\author{Truong Khang Tobias Nguyen}
\date{}

\begin{document}

\maketitle

\section{The Importance of Ethics}
I believe that ethics play an extremly important role for cybersecurity professionals. On that note, I also believe that as cybersecurity teams are amongst the most powerful entities in any organization even though they typically do not contribute to revenue generation. Elaborating on that, cybersecurity teams directly deal with the protections of an organizations digital assets which in most cases include information from individual people which the average person should not access.
\newline
Due to this, cybersecurity professionals receive a high level of trust as they are responsible to make sure that these assets can not be compromised. As they are tasked with the protection, cybersecurity professionals also enjoy a high degree of power over these assets. But here comes a dilemma - power enables control and surveillance. If desired, any cybersecurity professional can compromise assets themselves, control what measures an organization must take and even take down services themselves. The problem with that is that despite the power we give cybersecurity professionals, we also trust them to not abuse this power and use it against us. In a way, you can draw a parallel to weapon such as guns. Although you can use a gun to protect something, ultimately they are still capable of taking lives.
\newline
That being said, the dilemma of cybersecurity in ethics go further. Even though security is used as a mean to protect, protection always come with trade-offs. One security measure that can be taken is using employee monitoring software - even though one motivation to use such software is to prevent employees from clicking, accessing or providing things that they are not supposed to do but this software can also be used for surveillance such as tracking productivity. Another example is logging - the purpose of logging can be to trace and potentially recreating events but again, it can also be used for surveillance. Or penetration tests as they are used to test out how vulnerable systems actually are but they can also be used to find vulnerabilities that can then be exploited.
\newline
My point is that ethics are essential to cybersecurity as they explain the dynamics between security and power and subsequently make them understandable. What I want to say is that every security measure comes with trade-offs that one must consider when implementing them. An understanding of what assets such a measure protects, how it protects them, how this measure can be used to also harm these assets and accountability is what makes ethics so important in this field.

\section{My own values}
As someone with a background in software development, my values were always that I have a certain responsiblity towards the applications I develop. With that in mind, I always make sure that what I develop is robust, scalable, secure and that the code is readable for others.
\newline
In that sense, there are always trade-offs in development that I would have to consider such as prioritizing iteration speed over quality, feature development over security and so on. This is where one can draw a parallel what I have said in the previous section about trade-offs in security measures where one must understand the impact and consequences a measure can have.
\newline
That being said, to transition my values from the development perspective to the cybersecurity perspective, I believe in the following values:
\begin{itemize}
    \item \textbf{Protection:} Considering that protecting is the main objective of cybersecurity, I would of course put this into the list of cybersecurity values I have. Considering that others trust us with information and data they wouldn't give to others, requires us to respect the trust we receive and do everything in our hands to protect them and not betray the trust given to us.
    \item \textbf{Privacy:} Even though we hold the power of the assets we protect also gives us control over them, I believe that our task is to protect and not to pry on other's information. We should not try to leverage assets for purposes other than protection, such as surveillance, blackmailing or personal gain.
    \item \textbf{Fairness:} As the assets we protect are connected to multiple components of a systems, we have to ensure that our security measures treat all components equally without discriminating other components in favour of others. Furthermore, every security measure we implement should then keep in mind that not every component has the same level of knowledge for a measure and thus, it is also our responsibility to teach them.
    \item \textbf{Transparency:} In order for others to understand why we implement certain security measure, we need to be transparent in all the decisions and reasoning taken in the process. I believe that maintaining such transparency will result in those who we protect to accept the measures we take and will subsequently affect them.
    \item \textbf{Accountablity:} Another important value is accountability as it helps to identify the right instances answer for security outcomes rather then to dump the risk and blame onto one instance. Taking responsibility for our decisions and subsequent actions is essential, especially when mistakes or incidents happen. Rather than to use accountability to find the right instance to blame, it should be used to encourage learning to not let the same mistakes or incidents happen in the future.
    \item \textbf{Integrity:} I also believe that maintaining honesty and integrity when disclosing vulnerabilities, incidents and audits is important in maintaining the trust given to us. Even though disclosing weaknesses may reflect poorly on those we protect, I believe that it helps to address these problems properly and ensure that the necessary actions to alleviate them can be taken.
    \item \textbf{Ownership:} Considering that cybersecurity nowadays does not only include technical aspects. Human aspects are becoming more and more important in the cybersecurity landscape as even the most sophisticated technical defenses fall apart when human failure is involved. That being said, involving the people we protect rather than to see them as adversaries or relying on technical solutions will contribute to better cybersecurity measures and culture.
\end{itemize}
On that note, I realize that we already talked about some of these values in the lecture but these points resonate with me and thus, became my own values towards cybersecurity.
\newline
In conclusion, these are values important to me as I believe that they are essential in fostering a healthy cybersecurity environment that not only helps us to properly protect assets but to also foster collaboration with the ones we need to protect and improve together.

\end{document} 